The active experiments are run from the project notebooks and Python modules in
the repository.
The current implementation uses a compact extractive QA setup with a default
DistilBERT backbone and SQuAD~2.0 preprocessing.

\paragraph{Dataset split.}
The training pipeline creates deterministic train, validation, and development
splits from SQuAD~2.0.
For the answer-only baseline, the training set is restricted to answerable
examples.
For the abstain-aware setting, both answerable and unanswerable cases are kept.
For Stage 2, the verifier is trained on examples derived from the Stage 1 QA
outputs so that each instance contains a question, a passage, a predicted
answer, and a support label.
For Stage 3, the Stage 1 abstain-aware model and the Stage 2 verifier are kept
fixed, and calibration is fit post hoc on the validation split before raw and
calibrated confidence behavior is reported on the development split.

\paragraph{Evaluation.}
Model outputs are post-processed into final decisions and scored with:
\begin{itemize}[leftmargin=1.25em]
  \item answerable exact match,
  \item answerable F1,
  \item decision-aware overall exact match,
  \item decision-aware overall F1,
  \item unsupported answer rate, and
  \item abstain precision, recall, and F1.
\end{itemize}
The decision-aware overall metrics give full credit to correct abstentions on
unanswerable examples, so they do not penalize uncertainty by construction.
For Stage 2, we additionally report support accuracy, precision, recall, and
F1, together with supported-answer rate and contradiction rate for the verifier
and its downstream gated QA behavior.
For Stage 3, we additionally report expected calibration error (ECE), adaptive
ECE, maximum calibration error (MCE), Brier score, threshold-gap statistics,
score-label correlation, and reliability diagrams for both the QA decision
signal and the support signal.
For the next comparison stage, the same core metrics should be kept while
adding a controller-facing view:
supported-answer rate among answered cases, answer rate, and the selected
unsupported-answer budget.

\paragraph{Execution workflow.}
The current team workflow runs experiments from browser Google Colab while
editing source code locally in VS Code.
The notebook supports both hosted Colab and Colab connected to a local runtime;
the Stage 1, Stage 2, and Stage 3 runs reported here used a local runtime
backed by the project's WSL environment and GPU.
Artifacts are stored under the project runtime directory, and code changes are
synchronized through Git branches.

\paragraph{Next-stage comparison design.}
The next experimental comparison should hold the backbone, split, tokenizer,
and evaluation protocol fixed while varying only the decision formulation.
The most informative comparison is:
\begin{enumerate}[leftmargin=1.25em]
  \item the current modular pipeline
  (Stage 1 abstention, Stage 2 verifier, Stage 3 calibration),
  \item a fixed constrained controller over calibrated QA and support signals,
  and
  \item a jointly trained support-constrained multi-head learner that keeps
  the same encoder, tokenizer, split, and core QA supervision but replaces
  post-hoc control with an answer-support-abstain objective.
\end{enumerate}
Under this design, the key question is not whether one more module can be
bolted on, but whether a direct constrained objective improves the final
answer-support-abstain trade-off beyond what post-hoc modular tuning can
achieve.

\paragraph{What to lock before the final paper.}
Before turning this draft into a submission-ready paper, we should freeze:
\begin{itemize}[leftmargin=1.25em]
  \item the final model checkpoint choice,
  \item the exact QA, verifier, calibration, and controller-selection policies,
  \item the final metric tables and reliability figure,
  \item at least one qualitative error analysis slice, and
  \item the figure or table that best shows the utility-versus-groundedness
  trade-off for modular versus constrained-learning approaches.
\end{itemize}
